%======================================================================
\section{What v10 adds, and what it does not}
\label{sec:whatv10}
%======================================================================

v10 keeps every v9 statement \cite{CheungConeV9}. The v9 statements that v10 amends are
listed in Section~\ref{sec:v10changes}. The new items, with their
status, are:

\begin{center}
\small
\begin{longtable}{@{}L{.30\linewidth}L{.62\linewidth}@{}}
\caption{v10 verdicts (all executed; hashes in Appendix~\ref{app:v10artifacts}).}\label{tab:v10verdict}\\
\toprule
Item & Status\\
\midrule
\(\Delta J\) closed form & CR mode \emph{derived}. F mode \emph{identified} by PSLQ and verified at 9 fresh angles (\(\le1.5\times10^{-52}\)), not derived. Section~\ref{sec:v10dJ}.\\
\(\lambda_{\rm CB},\lambda_{\rm CK}\) & \emph{Derived} from tree operators and the v9 dictionaries (Theorem~\ref{thm:v10lambda}).\\
Exceptional points inside chambers & \emph{None}: rank 26 on all \(N=34\) chambers, 29 on F above \(\tfrac12\). Conditional on the a priori degree bound (Theorem~\ref{thm:v10chamber}).\\
Ninth cocycle on F above \(\tfrac12\) & \emph{Resolved}. It is present at every angle; it is not a constant vector, while the other eight are.\\
Transverse-jet catalogue & \PARTIAL. Exact slice jets for all 9 graphs. Same-channel \(r^2\) part dressed in two declared schemes; mixed and weight \(\ne0\) parts \OPEN{} (Section~\ref{sec:v10jet}).\\
D04 part of second \(\phi\phi\) value & \PARTIAL. The same-channel part is exactly 0. The full jet term is mixed \((1,0)\), with constant density \(\kappa^2/4\pi\). Per-mode traces are computed; the mode sum is not executed.\\
20 jet Ward combinations; 34-/18-dim maps; \(e_5,e_6\) & \OPEN. Precise obstruction in Section~\ref{sec:v10conseq}. No residual is claimed, since none was evaluated.\\
Vector-key finite parts & Computed in a declared \(k\)-\(\zeta\) scheme at 4 angles (Table~\ref{tab:v10vec}).\\
Registry & Unchanged (258 ids); no new sheet directory written.\\
\bottomrule
\end{longtable}
\end{center}

%======================================================================
\section{v10: closed form of the finite TT branch difference}
\label{sec:v10dJ}
%======================================================================

\begin{proposition}[CR mode; derived]
\label{prop:v10ICR}
Let \(\mu=1/\alpha\) and \(\beta=\mu-1\). Then
\[
 I(\mu,\mu-2)=\frac{(\mu-1)(6\mu^2+3\mu-2)}6-(\mu-1)^3(\mu+1)\,\psi_1(\mu-1).
\]
The same expression gives every \(k\le-2\) mode with \(\mu=m/\alpha\).
\end{proposition}

\begin{proof}[Derivation]
In the Gamma series of \loc{code/v9/modefun.py}, use
\((n+1)!\,\Gamma(\mu+1)/\Gamma(\mu+1+n)=\mu(\mu-1)B(n+2,\beta)\).
Each sum becomes a log-Beta integral. Continue it analytically,
keeping the Laurent data at the coincident poles. With
\(T=\psi_1(\beta)\), the intermediate sums are
\(G_2=(T\beta-1)/\beta\),
\(G_3=(-T\mu^2+3T\mu-2T+\mu-\tfrac32)/\beta\) and
\(G_4=(6T\mu^3-36T\mu^2+66T\mu-36T-6\mu^2+27\mu-25)/(12\beta)\).
Collecting them gives the formula
(\loc{code/v10/cr_mode_closed.py}). At \(\alpha=3/5\),
\(10000/7071\), \(4/7\) and \(99/100\) it matches the 34-digit
series to \(\le6\times10^{-35}\).
\end{proof}

\begin{proposition}[F mode; identified, not derived]
\label{prop:v10IF}
\[
 I(\mu,2-\mu)=\frac{(\mu-1)(59\mu^3-57\mu^2-62\mu+48)}{36}-(\mu-1)^3(\mu+1)\frac{\pi^2}6 .
\]
\end{proposition}
PSLQ \cite{PSLQ} on the basis \(\{I,1,\pi^2\}\) found integer relations at
\(\mu\in\{6/5,5/4,4/3,7/5,3/2,8/5,5/3\}\). The values were computed
to 55 digits by direct quadrature of Proposition~\ref{prop:v9g}
(\loc{code/v10/f_mode_pslq.py}). The relations interpolate to the
displayed polynomial. That polynomial was then checked at \(7/4\)
and \(9/5\), which were not used in the fit.
This is an identification, and its status is \emph{verified, not
derived}.

\begin{corollary}[\(\Delta J\) in closed form]
\label{cor:v10dJ}
\[
 \Delta J(\alpha)=\frac{(\mu-1)(-59\mu^3+93\mu^2+80\mu-60)}{36}
 -(\mu-1)^3(\mu+1)\Bigl[\psi_1(\mu-1)-\frac{\pi^2}6\Bigr],\qquad \mu=\tfrac1\alpha .
\]
It vanishes at \(\alpha=\tfrac12\) (\(\psi_1(1)=\pi^2/6\), and the
polynomial has a zero at \(\mu=2\)). As \(\alpha\to1\),
\(\Delta J=-\tfrac12(\mu-1)+O((\mu-1)^2)\), consistent with the v9
value \(\Delta J(99/100)\approx-4.9\times10^{-3}\).
\end{corollary}

\begin{center}\scriptsize
\begin{longtable}{@{}llll@{}}
\caption{Fresh-angle verification (\loc{code/v10/dJ_closed_verify.py}; quadrature at 60 digits, and 80 digits for a precision check).}\label{tab:v10dJ}\\
\toprule
\(\mu=1/\alpha\) & \(|I_{\rm F}^{\rm quad}-I_{\rm F}^{\rm closed}|\) & \(|I_{\rm CR}^{\rm quad}-I_{\rm CR}^{\rm closed}|\) & \(\Delta J\) (closed)\\
\midrule
@@DJFRESH@@
\bottomrule
\end{longtable}
\end{center}
The 14 v9 values of Table~\ref{tab:v9dJ} agree with the closed form
to \(4.4\times10^{-27}\). The v9 sentence ``we claim no elementary
closed form'' is therefore superseded for the CR mode and qualified
for the F mode.

%======================================================================
\section{v10: the ghost normalizations from tree data}
\label{sec:v10lambda}
%======================================================================

\begin{theorem}[Derivation of \(\lambda_{\rm CK}\) and \(\lambda_{\rm CB}\)]
\label{thm:v10lambda}
Use the dictionaries of Sections~\ref{sec:v8map} and
\ref{sec:v9ghost}, with signs \(s_i=1\), \(s_c=-1\), \(s_{\bar c}=1\);
the \(K\)-dictionary \(K^{00}=-\mathrm{LL}_{\tau\tau}\),
\(K^{03}=-2\ii\,\mathrm{LL}_{\tau z}\),
\(K^{33}=\mathrm{LL}_{zz}\); and \(q_0=\ii q_\tau\), \(q_3=q_z\).
Then the following three tree-level operator identities hold on all
\(m=0\) atoms and all nine test momenta \((q_\tau,q_z)\in\{1,2,3\}^2\):
\begin{align*}
 \text{engine }K_2c&=\rho\,\delta_ch,& \delta_ch_{mn}&=p_mc_n+p_nc_m, & \rho&=-\ii,\\
 \text{engine }E(-q)h&=\epsilon\,(2-\delta_{mn})E^{mn}_{\rm cat}(h), && & \epsilon&=1,\\
 \text{engine }C_2^{\mathsf T}(-q)b&=\tau\,Y_b, & Y_b&=-2\eta F_L^{\mathsf T}b, & \tau&=-\ii/2 .
\end{align*}
With the eval normalizations of the engine (\(\Gamma_{\rm sym}=4\Gamma\)
on HH and \(2\Gamma\) on CB and CK), the \(hc\) row becomes
\(8\rho B+2\lambda_{\rm CK}\epsilon G=0\). Theorem~\ref{thm:ST}
gives \(B+\lambda_{\rm pred}G=0\), with
\(\lambda_{\rm pred}=N/(\ii A_{\rm GH})=-2\ii\)
(Input~\ref{input:tree}). Hence
\[
 \lambda_{\rm CK}=\frac{4\rho\lambda_{\rm pred}}{\epsilon}=-8 .
\]
The \(b c\) row becomes
\(2\lambda_{\rm CB}\Gamma_{\rm D06}+2\lambda_{\rm CK}\tau\Gamma_{\rm D06}=0\)
by the ghost identity
\(\hat B_{\rm D06}=-2\eta F_L\hat\Gamma_{\rm D13}\)
(Theorem~\ref{thm:ghost}). Hence
\[
 \lambda_{\rm CB}=-\tau\lambda_{\rm CK}=-4\ii .
\]
\end{theorem}

\begin{proof}
\loc{code/v10/lambda_derive.py} evaluates the engine operators of
\loc{code/v8/apex_exact.py} (\(K_2\) at \(q\); \(E\) and
\(C_2^{\mathsf T}\) at \(-q\), as in the row construction of
\loc{apex_certificate.py}). Each ratio is single-valued on every
nonzero component, with no mismatch
(\loc{logs/v10/lambda_derive.json}). The two row reductions are
linear algebra on these ratios.
\end{proof}

\begin{remark}
The derivation uses no loop value. It does use the dictionary
signs, which v9 selected by exact solvability on loop data. The
values \(-8\) and \(-4\ii\) are therefore no longer free fitted
parameters, but they are not independent of the v9 sign choice.
The relation \(\lambda_{\rm CK}=\lambda_{\rm pred}\lambda_{\rm CB}\)
noted in v9 is explained by the numerical coincidence
\(4\rho/\epsilon=-4\ii=\lambda_{\rm CB}\), i.e.\ \(-\tau\lambda_{\rm CK}=4\rho/\epsilon\) with
\(\lambda_{\rm CK}=(4\rho/\epsilon)\lambda_{\rm pred}\).
\end{remark}

%======================================================================
\section{v10: chamber ranks without exceptional points}
\label{sec:v10chamber}
%======================================================================

\loc{code/v10/minor_poly.py} picks, at one angle per chamber, three
nonsingular maximal minors with \(\alpha\)-free row tags. The atom
key is \((C,n,j,i)\) with \(j=(m-\varepsilon|a|)/2\). The program
then re-evaluates the minors at 11--15 angles per chamber
(\loc{logs/v10/minor/<tag>/}).

\begin{theorem}[Rank inside chambers]
\label{thm:v10chamber}
Assume the a priori bound (D): every ST row entry is a polynomial of
degree \(\le2\) in \(1/\alpha\) on each chamber. This is argued from
the construction: \(E\) carries two derivative factors, \(K_2\) one,
and each is linear in \(1/\alpha\). It is not separately
machine-checked. Then:
\begin{enumerate}[label=(\alph*)]
\item On CR \((\tfrac14,\tfrac13),(\tfrac13,\tfrac12),(\tfrac12,\tfrac23),(\tfrac23,1)\)
and on F \((\tfrac14,\tfrac13),(\tfrac13,\tfrac12)\) (\(N=34\)), all
three minors are the same constants at every sampled angle,
\(124515522497539473408\,\ii\),
\(-73477495818488684950134624092160000\) and
\(-4059104145057262382496562859212800\), and their entries are
\(\alpha\)-independent. The rank is therefore 26 at \emph{every}
\(\alpha\) in these chambers.
\item On F \((\tfrac12,\tfrac23)\) and \((\tfrac23,1)\) (\(N=38\)),
minor~2 equals \(-400569487999071945641108176896000000000\,\ii\) at
all 28 sampled angles. Its only \(\alpha\)-dependent entries lie in
one column, so it is affine in that column and constant by (D). Hence
the rank is \(\ge29\) everywhere. Minors 0 and 1 vary, with 13 and
15 distinct values, and are never zero.
\item On F above \(\tfrac12\), \(\dim Z(\alpha)=9\) at 16 angles
(\loc{code/v10/cocycle9.py}). Every entry fits an affine function of
\(1/\alpha\) on 8 points per chamber. Write the 9 Cramer kernel
vectors of minor~2 at \(t=1/\alpha\). The residual \(W(t)\) of the
remaining rows has degree \(\le6\) by (D) and vanishes at 8 points,
so it vanishes identically. Hence the generic rank is 29, and the
rank equals 29 at every \(\alpha\): at most 29 by semicontinuity,
at least 29 by (b).
\item The intersection \(\bigcap_\alpha Z(\alpha)\) over the 8
angles of each chamber is 8-dimensional. Eight cocycles are constant
vectors. The ninth exists at every angle, but its direction depends
on \(\alpha\) and is rational in \(\alpha\) by Cramer's rule.
\end{enumerate}
\end{theorem}
This replaces the v9 range ``between 29 and 30'' by exactly 29,
conditional on (D). It also answers the v9 open item ``exceptional
set inside chambers'': the set is empty, conditional on (D).

%======================================================================
\section{v10: the transverse-jet extension of the slice catalogue}
\label{sec:v10jet}
%======================================================================

\paragraph{Construction.}
\loc{code/v10/jet_replay.py} executes the numerator section of the
frozen v7/v8 replay with one change. The external transverse momenta
\(p_1,p_2\) are kept as symbols instead of being set to zero. The
bubble denominator gains \(-p_\perp^2\), and the loop momentum stays
on the longitudinal slice (\(l_\perp=0\)). This is a declared
coverage limitation, and the status remains
\textsf{ALGEBRAIC\_SLICE\_ONLY}. The UV pole is then an exact
polynomial in \(q_0,q_3,p_\pm\). It is stored per key and channel
pair, by transverse degree and helicity weight
(\loc{logs/v10/jet/jet_<G>.json}). No frozen file is edited.

\paragraph{Checks.}
(i) At order \(p_\perp^0\), \(\sum_s\mathrm{pole}^{(0)}_{ss}z_s\)
reproduces the v8 catalogue on every computed key: 684 key--branch
checks with 0 mismatches. (ii) In D04 the UNIT sum is
\(-\kappa^2(q_0^2-q_3^2-p_\perp^2)/4\pi\), which is flat
covariance restored by the jet terms.

\paragraph{Census.}
Summed over the 9 graphs, the jet entries are:
\begin{itemize}
\item same-channel \((2,0)\): 128;
\item mixed \((2,0)\): 84;
\item same \((2,\pm2)\): 152 each;
\item mixed \((2,\pm2)\): 70 each;
\item same \((1,\pm1)\): 252 each;
\item mixed \((1,\pm1)\): 206 each.
\end{itemize}
D12 has none.

\begin{definition}[Declared jet dressing and schemes]
A same-channel, weight-0 coefficient of \(p_+p_-\) on channel \(s\)
is dressed by the \(r^2\) moment \(d_s\). In scheme A (the Abel
subtraction of Theorem~\ref{thm:D}),
\(d_0=\Dcal_0\), \(d_1=\Dcal_\alpha\),
\(d_2^{\rm F}=\Dcal_{2\alpha-1}\) and
\(d_2^{\rm CR}=\Dcal_{2\alpha}\). In scheme Z, each mode's
\(\Tr[r^{2w}e^{-sA_\mu}]\) is continued from \(\Re w<-\tfrac12\) to
\(w=1\), and the modes are summed by Hurwitz zeta. This mirrors the
role of \(\Ccal_b\) at constant profile. It is a reduction contract
in the sense of Assumption~\ref{ass:reduction}, not a theorem.
\end{definition}

\begin{proposition}[Scheme Z; executed]
\label{prop:v10Z}
Per mode, \(\Tr[r^2e^{-sA_\mu}]^Z=s\,\mu(\mu^2-1)/3\), from the
Mellin closed form
\((2s)^w\Gamma(w+1+\mu)\Gamma(-\tfrac12-w)/(2^{w+2}\sqrt\pi\,\Gamma(\mu-w))\).
This form matches quadrature for \(w<-\tfrac12\) to \(10^{-17}\) or
better. Its only poles are at half-integer \(w\), so \(w=1\) is
regular. The mode sum gives
\[
 \Dcal^Z_b=\Dcal_b+\frac{11\alpha}{180}\quad\text{(exactly, all }b),
\]
so the scheme dependence is a bulk term linear in \(\alpha\). Flipping one mode
\(\nu\mapsto-\nu\) reproduces \(\Delta\Dcal=2\nu(1-\nu^2)/3\) of
Proposition~\ref{prop:Dflip}.
\end{proposition}

\begin{proposition}[Dressed same-channel jets; executed]
\label{prop:v10dressed}
In scheme A, all dressed same-channel values vanish at \(\alpha=1\)
(F), and CR\(=\)F at \(\alpha=\tfrac12\), on every key
(\loc{logs/v10/jet/jet_analysis.json}). Examples:
D05a \((0,0)\) gives
\(\kappa^2(\alpha-1)(2\alpha-1)/(6\pi\alpha)\) (CR) and
\(\kappa^2(\alpha-1)(2\alpha-1)(\alpha^2-3\alpha+1)/(6\pi\alpha^3)\) (F),
each times \(p_+p_-\). D11 \((1,1)\) gives
\(\kappa^2(\alpha^2-1)(11\alpha^2+1)/(2880\pi\alpha^3)\) in scheme A
and \(-\kappa^2(10\alpha^2+1)/(2880\pi\alpha^3)\) in scheme Z. Of the
74 keys with weight-0 \(r^2\) content (CR), 38 are complete. The
other 36 (all of D01, D06, and 9 of D03) also need mixed pairs and
are \PARTIAL.
\end{proposition}
The D05a value is a slice-jet coefficient under the declared
dressing. It is not compared with the earlier retained statement
``D05a part \(=0\)''. That statement came from a different reduced
model (a large-radius count without the slice), and no
supersession is claimed.

\paragraph{The mixed \(r^2\) trace.}
Term-by-term integration gives
\[
 T_1(x;\mu,\nu)=\tfrac14(4x(1-x))^2\tfrac{\Gamma(h+2)\Gamma(h+1)}{\Gamma(\mu+1)\Gamma(\nu+1)}x^\mu(1-x)^\nu
 F_4(h+2,h+1;\mu+1,\nu+1;x^2,(1-x)^2)
\]
with \(h=(\mu+\nu)/2\). The double series has power divergences
\(N^{1/2-j}\) and no logarithms. Write \(F_4(h+2,h+1)=(h+1+\theta)(h+\theta)F_4(h,h+1)/(h(h+1))\), with \(\theta\) the
Euler operator, and apply Bailey's reduction \cite{BaileyF4} off the boundary,
\(\lambda=1-t^2\). The finite part is then
\(h(h+1)\phi_0-2(h+1)\phi_2+2\phi_4\) in the Taylor coefficients
of \(\Phi(t)\) (\loc{code/v10/t1_fp.py}). It agrees with an
independent shell-by-shell zeta continuation to
\(10^{-19}\)--\(10^{-24}\) (\loc{logs/v10/jet/t1_fp_check.json}). It
reproduces \(\mu(\mu^2-1)/3\) exactly at \(\mu=\nu\) for every
\(x\), and its endpoint values are \(T_1(0)=\mu(\mu^2-1)/3\) and
\(T_1(1)=\nu(\nu^2-1)/3\). After \(x\)-integration, the shell
continuation is \emph{not} reliable, because of non-uniform
endpoint asymptotics (\loc{logs/v10/jet/mixed_r2.json}). It is used
nowhere.

%======================================================================
\section{v10: consequences of the jets}
\label{sec:v10conseq}
%======================================================================

\paragraph{D04.}
The constant density for the jet term is
\(\rho_{(1,0)}=\kappa^2/4\pi\) (\loc{code/v10/jet_density_D04.py}).
Under the declared rule, the D04 part of the second \(\phi\phi\)
value is
\(\tfrac{\kappa^2}{4\pi}M_{10}\), where
\(M_{10}=\sum_{k\ge0}\varphi(k/\alpha)+\sum_{j\ge1}\varphi(j/\alpha-1)\)
is \(k\)-\(\zeta\)-regulated and \(\varphi(n)=\int_0^1T_1(x;n+1,n)\,\mathrm dx\).
It is branch independent, since only channels 0 and 1 enter. We
have \(\varphi\) at 27 points (for example
\(\varphi(0)=-0.0578297132876329558555\) and
\(\varphi(3/2)=2.32785224733468215021\);
\loc{logs/v10/jet/phi/}). No closed form was found in the basis
\(\{n^j,n^j\psi_1(n),n^j\psi(n)\}\) (\loc{code/v10/phi_fit.py}). The
large-\(n\) runs needed for the mode sum were stopped for time.
\(M_{10}\) is therefore \emph{not evaluated}, and the D04 part is
\PARTIAL. A further caveat: the bare \(r^2\) insertion between
channels 1 and 0 omits the frame intertwining \(H_I\), which the
spin-weight-one channel may require. The candidate is therefore a
declared-rule quantity, not a derived value.

\paragraph{The 20 jet-level Ward combinations, the 34- and 18-dimensional
projections, and \(e_5,e_6\).}
These need catalogue values on the engine's \(m=2\) and weight
\(\pm1\) atoms. The jet catalogue now supplies the weight-0
same-channel part (Proposition~\ref{prop:v10dressed}). It does not
supply the following:
\begin{enumerate}[label=(\alph*)]
\item the mixed weight-0 pairs (84 entries), which need \(M_{st}\)
mode sums as for D04;
\item weight \(\pm1\) (412 mixed, 504 same entries). These need
\(k\to k\pm1\) traces with an \(r\,e^{\pm\ii\theta}\) insertion, at
\(w=\tfrac12\), where scheme Z has a pole: a genuine
logarithm;
\item weight \(\pm2\), which needs \(k\to k\pm2\) traces;
\item a dictionary from jet keys to engine \(m=2\) atoms. Such a
dictionary was not constructed.
\end{enumerate}
No residual was evaluated, and none is reported as zero. Status:
\OPEN, with the obstructions (a)--(d).

%======================================================================
\section{v10: vector-key finite parts in a declared scheme}
\label{sec:v10vector}
%======================================================================

\begin{proposition}[Executed]
\label{prop:v10vector}
Per mode, take scheme Z with \(\Tr_{\rm mode}e^{-A_\mu}=-\mu/2\); this
is consistent with \(\int g=1/(2(\mu+\nu))\). Then for a pair density
\(\rho\) with \(R=\int_0^x\rho\),
\(I_\rho=-R(1)\nu/2-(\mu^2-\nu^2)\int_0^1Rg\). Pair the modes at
the same \(k\) (channel 0: \(|k|/\alpha\); channel 1:
\(|k+\alpha|/\alpha\)), and regulate the sum by \(|k|^{-\varepsilon}\)
(\loc{code/v10/vector_fp.py}, \loc{vector_combine.py}). Then, for
D01\(+\)D03 in the v7 normalization:
\begin{enumerate}[label=(\alph*)]
\item the \(1/\varepsilon\) residues equal the v9 logarithms
\((\alpha/2)L_{\rm key}\) of Proposition~\ref{prop:v9vector} to
\(\le1.1\times10^{-20}\) at all four angles. For keys \((1,1)\) and
\((2,2)\) the pole is \(-\alpha\kappa^2P_L^2/(60\pi\varepsilon)\);
\item the off-diagonal keys \((1,6),(2,8),(6,1),(8,2)\) have finite
part 0, because D01 and D03 cancel digit for digit;
\item the diagonal keys have finite part
\(+c(\alpha)\kappa^2P_L^2\) on \((1,1),(2,2)\) and
\(-c(\alpha)\kappa^2P_L^2\) on \((6,6),(8,8)\), with \(c\) given in
Table~\ref{tab:v10vec}. They are branch independent.
\end{enumerate}
\end{proposition}
\begin{table}[ht]\centering\small
\caption{Vector-key finite parts (declared \(k\)-\(\zeta\) scheme).}\label{tab:v10vec}
\begin{tabular}{@{}ll@{}}
\toprule
\(\alpha\) & \(c(\alpha)\) (\(k\)-\(\zeta\) finite part)\\
\midrule
@@VECTAB@@
\bottomrule
\end{tabular}
\end{table}
The finite part depends on the scheme. Rescaling the regulator
\(|k|\to\lambda|k|\) shifts \(c\) by
\((\alpha/60\pi)\log\lambda\), and a sharp cutoff differs from it
by the \(\zeta(-1),\zeta(0)\) terms. It is a declared-scheme value
and not an apex coefficient.

%======================================================================
\section{v10: changes to v9 statements}
\label{sec:v10changes}
%======================================================================
\begin{itemize}
\item Table~\ref{tab:v9dJ} discussion: ``PSLQ was not attempted;
no elementary closed form'' is superseded by
Corollary~\ref{cor:v10dJ} (CR derived, F identified).
\item Theorem~\ref{thm:v9ghost}: ``two fitted normalizations'' is
now ``two normalizations derived from tree operators''
(Theorem~\ref{thm:v10lambda}).
\item Theorem~\ref{thm:v9generic}: ``F\((>1/2)\) between 29 and 30''
is now ``exactly 29 at every angle'', conditional on (D)
(Theorem~\ref{thm:v10chamber}).
\item Section~\ref{sec:v9D04}: obstruction (a), ``the slice does not
export jet densities'', is removed by the jet replay. Obstruction
(b) is made precise: the D04 jet term is purely mixed, and scheme Z
regularizes the per-mode power divergence. The mode sum is still
missing.
\item Proposition~\ref{prop:v9vector}: the finite parts are now
computed in a declared scheme.
\end{itemize}
No v9 number was found to be in error. The frozen-tree digests
before and after v10 are in Appendix~\ref{app:v10artifacts}.
